\section{Preliminaries}
\label{sec:preliminary}

We work over a prime-order Galois field $\Field$,
with multiplicative group $\Field^*$.
(Row) vectors are denoted by bold letters,
\eg, $\bm{a}, \bm{b}, \bm{c}$,
and scalars by lowercase letters, \eg, $a, b, c$.
We write $[n] := \{1, \ldots, n\}$.
% The inner product operation is denoted by $\langle\cdot,\cdot\rangle$.

Given a point $\bm{x}\in \Field^\mu$,
we derive a $2^\mu$-dimensional tensor-structured vector
via the following reverse tensor span.

\begin{definition}[Reverse Tensor Span]
	\label{def:tensorspan}
	Given $\mu$ values $\bm{x} \in \Field^\mu$,
	their \textbf{reverse tensor span} is the vector
	$\tensorspan{\bm{x}} \in \Field^{2^\mu}$ defined as:
	\[
		\tensorspan{\bm{x}}
		:= (1, x_\mu) \kronecker (1, x_{\mu - 1}) \kronecker \cdots \kronecker (1, x_1).
	\]
\end{definition}

Given a vector $\bm{a} \in \Field^n$, let
$q:=\mathsf{next\_power\_of\_two}(n)$.  For a point
$\bm{x} \in \Field^{\log_2 q}$, we define the \emph{evaluation} of $\bm{a}$ at
$\bm{x}$ as an inner product with its conceptual zero extension:
\begin{equation}
\label{eq:evaluation}
    \tensorproj{\bm{a}}{\bm{x}}
    := \left\langle \bm{a}\| \tensorspan{\bm{x}},\bm{0}^{q-n} \right\rangle.
\end{equation}

\mysmallskip
\noindent
\textbf{MLE and multilinear PCS.}
There is a one-to-one correspondence between vectors (or matrices)
and multilinear polynomials.
The multilinear extension (MLE) of a vector $\bm{a}$
is denoted by $\psi_{\bm{a}}(\bm{x})$.

\begin{definition}[MLE~\cite{crypto/Thaler13}]
The MLE $\psi_{\bm{a}}(\bm{x})$ of a vector $\bm{a} \in \Field^{n}$, where
$q:=\mathsf{next\_power\_of\_two}(n)$, is the unique multilinear polynomial
$\Field^{\log_2 q} \to \Field$
such that:
\[
	\psi_{\bm{a}} (\bm{i}) =
    \begin{cases}a_i,&i\le n,\\0,&n<i\le q,\end{cases}
	\quad \forall\, \bm{i} \in \{0, 1\}^{\log_2 q},
\]
where $\bm{i}$ ranges over the Boolean hypercube
and corresponds to the bit decomposition of the index $i$.
\end{definition}

Existing work~\cite{tcc/Lee21,asiacrypt/BuenzMMTV21,tcc/BootleCG20,eurocrypt/ChiesaHMVW20,crypto/Setty20}
shows that MLE evaluation can be expressed as an inner product.
Using the above definition of $\Eval(\bm{a}, \bm{x})$,
we have the following relation.

\begin{lemma}
\label{lemma:mleProjection}
	The MLE $\psi_{\bm{a}}(\cdot)$ of a vector $\bm{a} \in \Field^{n}$ satisfies,
    for $q:=\mathsf{next\_power\_of\_two}(n)$:
	\begin{equation}
	\label{eq:mleprojection}
    \psi_{\bm{a}}(\bm{x})
    =
    \Bigl[\prod_{i=1}^{\log_2 q}(1-x_i)\Bigr]
    \cdot \Eval(\bm{a}, \bm{x}).
	\end{equation}
\end{lemma}

Throughout the remainder of this paper, we ignore the scaling factor and
use $\Eval(\bm{a}, \bm{x})$ to indicate vector evaluation, without referring to the MLE $\psi_{\bm{a}}(\cdot)$.

Due to the correspondence between vectors and polynomials,
a PCS can also be viewed as a vector commitment scheme.
We define PCS using equivalent vector notation.

\begin{definition}[Multilinear Polynomial Commitment Scheme]
\label{def:pcs}
	A multilinear PCS for vector commitments
	is a tuple of algorithms:
	\begin{itemize}[leftmargin=*, labelsep = 0.5em]
	\item $\pp \leftarrow \Setup(1^\lambda, n)$
	takes as input the vector length $n$
	and outputs public parameters $\pp$;
	\item $C \leftarrow \Com(\pp, \bm{a})$
	(or $\Com(\bm{a})$ for short)
	takes as input a vector $\bm{a} \in \Field^{n}$
	and outputs a commitment $C \in \ComSpace$;
	\item $\mathsf{bit} \leftarrow \Open(\pp, C, v, \bm{x} : \bm{a})$
	is an interactive protocol between a prover $\Prover$
	and a verifier $\Verifier$,
	where $\Prover$ convinces $\Verifier$
	that $\Eval(\bm{a},\bm{x}) = v$.
    At the end of $\Open$, $\Verifier$ outputs $\mathsf{bit} \in \{ \TRUE, \FALSE \}$.
	\end{itemize}
\end{definition}

For simplicity, we write $\Com(\pp, \bm{a})$ as $\Com(\bm{a})$.

As a matrix can be flattened into a vector,
a PCS also serves as a matrix commitment scheme.

\mysmallskip
\noindent
\textbf{Inner product argument.}
Given committed vectors, an inner product argument provides proofs for their inner products.

\begin{definition}[Inner Product Argument]
\label{def:ipa}
Given a vector commitment scheme $\Com$
with public parameters $\pp$
and commitment space $\mathbb{C}$,
an IPA is a zkSNARK for the relation:
\begin{align*}
\mathcal{R}_{\sf ip} = & \left\{ \left( \begin{aligned} & \pp; c \in \Field, C_a, C_b \in \mathbb{C} \\ \end{aligned} \right) \right. : \left( \begin{aligned} & \bm{a}, \bm{b} \in \Field^n \\ \end{aligned} \right) \\
& \quad\left|~
\begin{aligned} c = \bm{a} \bullet \bm{b} \wedge C_a = \Com(\vec{\bm{a}}) \wedge C_b = \Com(\bm{b}) \end{aligned} \right\}.
\end{align*}
\end{definition}

An IPA can be instantiated by combining
an inner product proof with a PCS
(see Appendix~\ref{appx:litebullet}).

\mysmallskip
\noindent
\textbf{Collection commitment for heterogeneously sized vectors.}
Given a collection of heterogeneously sized vectors
$\{\bm{a}_j\}_{j\in [T]}$, let $n_j$ be the exact length and let $q_j$ be a
public power-of-two cover satisfying $n_j\le q_j<2n_j$.  We define a
collection-commitment interface that commits to all vectors jointly and
supports repeated batched openings at points $\{\bm{z}_j\}_{j\in[T]}$.  The
packed entries, aligned offsets, and exact matrix dimensions are authenticated
separately.  Applications may disclose the dimensions, or keep them committed
and reveal only the cover sizes.  With the hiding base commitment in
Appendix~\ref{appx:pcs} and zero-knowledge openings, the latter mode hides both
parameter entries and exact matrix sizes.

Additionally,
the prover commits to the evaluation values and points
via commitments $C_V$ and $C_X$,
so that the verifier only needs these commitments
to verify the openings.

\begin{definition}[Heterogeneous-Vector Collection Commitment]
\label{def:collectioncommit}
A collection commitment with collective MLE openings consists of the following
algorithms:
\begin{itemize}[leftmargin=*, labelsep=0.5em]
\item $\pp \leftarrow \Setup(1^\lambda, K)$
takes as input the public cover-domain capacity $K$
and outputs public parameters $\pp$.
\item $(C_a,C_I,C_D,\pi_{\mathsf{size}}) \leftarrow
\Com(\pp,\{q_j\}_{j\in[T]};\{r_j,c_j,n_j,\bm{a}_j\}_{j\in[T]})$
takes as public input the cover sizes and as private input the exact matrix
dimensions, lengths $n_j=r_jc_j$, and entries
$\bm{a}_j\in\Field^{n_j}$;
outputs a commitment $C_a$ that collectively commits to all packed vector
entries, a commitment $C_I$ to their aligned bucket offsets, a hiding
commitment $C_D$ to the exact dimensions, and a one-time zero-knowledge proof
$\pi_{\mathsf{size}}$ of their consistency with the public covers.
\item $C_X \leftarrow \mathsf{CommPoint}(\pp, \{\bm{x}_j\}_{j\in[T]})$
(or $\mathsf{CommPoint}(\{\bm{x}_j\}_{j\in[T]})$ for short)
takes as input the evaluation points $\{\bm{x}_j\}_{j\in[T]}$
and outputs a commitment $C_X$ that collectively commits to these points.
\item $C_V \leftarrow \mathsf{CommValue}(\pp, \{v_j\}_{j\in[T]})$
(or $\mathsf{CommValue}(\{v_j\}_{j\in[T]})$ for short)
takes as input the claimed evaluation values $\{v_j\}_{j\in[T]}$
and outputs a commitment $C_V$ that commits to these values.
\item $\mathsf{bit} \leftarrow
\Open(\pp, C_a, C_I, C_D, C_V, C_X:
\{\bm{a}_j\}_{j\in[T]},
\{v_j\}_{j\in[T]},
\{\bm{x}_j\}_{j\in[T]})$
is a protocol between a probabilistic-polynomial-time prover $\Prover$
and verifier $\Verifier$.
Both $\Prover$ and $\Verifier$ hold the commitments
$C_a, C_I, C_D, C_V, C_X$ and the cached proof $\pi_{\mathsf{size}}$.
The prover additionally knows
$\{r_j,c_j,n_j,\bm{a}_j\}_{j\in[T]}$, $\{v_j\}_{j\in[T]}$, and
$\{\bm{x}_j\}_{j\in[T]}$,
and aims to convince the verifier that
 $\Eval(\bm{a}_j, \bm{x}_j) = v_j$ for all $j \in [T]$.
At the end of $\Open$, the verifier outputs
$\mathsf{bit} \in \{\TRUE, \FALSE\}$.
\end{itemize}
\end{definition}

Our goal is to realize this interface so that physical prover work and memory
track the input entries rather than a materialized padded domain.

\mysmallskip
\noindent
\textbf{zkSNARKs and related notions.}
Standard definitions of arguments of knowledge,
zkSNARKs,
perfect completeness,
knowledge soundness,
zero knowledge,
and the Fiat--Shamir transformation
are recalled in Appendix~\ref{appx:definitions}.
